% TOP_PROBLEM: {"id":30007588,"problem_number":"LOCAL-30007588","title":"Laboratory-to-field transfer","category_id":21,"category":{"id":21,"name":"economics","display_name":"Economics","description":"Open economic theory statements, published research agendas, and proposed scoped specifications; question provenance and historical priority are qualified per record.","slug":"economics","order_index":21,"created_at":"2026-10-08T00:00:00Z"},"status":"research_question","external_url":null,"legacy_ids":[],"published":false,"statement":"Determine when laboratory treatment effects transport to field populations and environments, and estimate the transport gap under explicit selection assumptions.","economics":{"original_id":77,"field":"Behavioral and experimental economics","kind":"identification","formalization_basis":"adapted_research_question","source_status":"research_agenda","priority_status":"earliest_found","priority_note":"This 2007 primary paper predates the catalogue source and explicitly frames the question. Its acknowledgement mentions a 2006 precursor, so original priority is not claimed.","earliest_verified_reference":{"authors":["Steven D. Levitt","John A. List"],"title":"What Do Laboratory Experiments Measuring Social Preferences Reveal About the Real World?","year":2007,"url":"https://bpb-us-w2.wpmucdn.com/voices.uchicago.edu/dist/f/1276/files/2018/09/27-101-25264884-1gx76fi.pdf","doi":"10.1257/jep.21.2.153","locator":"Introduction, pp. 153–154; Stakes, p. 165; Conclusions, pp. 170–171 (PDF pp. 1–2, 13, 18–19).","evidence_summary":"The paper frames generalizability as a central question, identifies scrutiny, context, selection, moral considerations, and stakes as moderators, and calls for theory connecting laboratory and naturally occurring behavior."},"current_reference":{"authors":["Steven D. Levitt","John A. List"],"title":"What Do Laboratory Experiments Measuring Social Preferences Reveal About the Real World?","year":2007,"url":"https://bpb-us-w2.wpmucdn.com/voices.uchicago.edu/dist/f/1276/files/2018/09/27-101-25264884-1gx76fi.pdf","doi":"10.1257/jep.21.2.153","locator":"Introduction, pp. 153–154; Stakes, p. 165; Conclusions, pp. 170–171 (PDF pp. 1–2, 13, 18–19).","evidence_summary":"The paper frames generalizability as a central question, identifies scrutiny, context, selection, moral considerations, and stakes as moderators, and calls for theory connecting laboratory and naturally occurring behavior."},"formalization_note":"The source explicitly supplies a generalizability agenda. The potential-outcome transport formula is our operational formulation and requires assumptions not inferred from the paper’s existence."},"statusQualification":"Published question or research agenda verified at the cited locator. The mathematical specification is adapted for this catalogue and includes additional assumptions; source publication does not establish first-ever priority or certify this exact model remains unsolved. The source explicitly supplies a generalizability agenda. The potential-outcome transport formula is our operational formulation and requires assumptions not inferred from the paper’s existence.","statusReviewedAt":"2026-10-08"}
% FIRST_POSTED: 2026-10-08
% ENTRY_KIND: statement_only
% !TeX program = lualatex
\documentclass[11pt]{article}
\usepackage{fontspec}
\usepackage[a4paper,margin=25mm]{geometry}
\usepackage{amsmath,amssymb,mathtools}
\usepackage{xurl}
\usepackage[hidelinks]{hyperref}
\newcommand{\sref}[2]{\hyperlink{source:#1:#2}{[S#2]}}
\setlength{\emergencystretch}{3em}
\begin{document}
% problemId:30007588
\section*{Laboratory-to-field transfer}\label{30007588}
\subsection{Definitions and mathematical statement}
Let \(S\in\{L,F\}\) indicate a laboratory sample or a target field population. For a specified intervention \(D\in\{0,1\}\), let \(Y_i(d,S)\) denote individual \(i\)'s potential outcome under treatment \(d\) in setting \(S\). Let \(X_i\) collect participant characteristics and \(C_S\) collect contextual features, including scrutiny, stakes, framing, and institutional rules. The field target and a candidate laboratory transport formula are
\[\tau_F=\mathbb E_F[Y_i(1,F)-Y_i(0,F)],\qquad \widetilde\tau_F=\int\mathbb E[Y_i(1,L)-Y_i(0,L)\mid X_i=x]\,dP_F(x).\]
The formula requires overlap in \(X\), internal identification of the laboratory contrast, and invariance of conditional treatment effects across contexts. It therefore does not follow merely from random assignment inside the laboratory. Determine which observed contextual moderators permit a credible extended transport model \(\tau(x,C)\), and quantify errors when context or selection variables are unmeasured. Use deliberately varied laboratory settings and independent field comparisons to evaluate transported predictions, with the treatment and outcome kept conceptually comparable. Distinguish differences in effect magnitudes from different outcome measurement or different versions of the intervention. The research agenda is to identify defensible domains of extrapolation and sensitivity bounds for specific behaviors, including social-preference games, rather than to prove or deny external validity of laboratory economics as a whole.

\par\textbf{Formulation and scope.} The source explicitly supplies a generalizability agenda. The potential-outcome transport formula is our operational formulation and requires assumptions not inferred from the paper’s existence.

\par\textbf{Open-question provenance.} An explicit question or research agenda was verified at the source locator. This does not by itself certify that every later version remains unresolved \sref{30007588}{1}.

\par\textbf{Historical priority.} This 2007 primary paper predates the catalogue source and explicitly frames the question. Its acknowledgement mentions a 2006 precursor, so original priority is not claimed.

\subsection{Short English statement}
Determine when laboratory treatment effects transport to field populations and environments, and estimate the transport gap under explicit selection assumptions.

\subsection{Sources}
\begin{itemize}\raggedright
\item[\textnormal{[S1]}]\hypertarget{source:30007588:1}{} Steven D. Levitt, John A. List. 2007. What Do Laboratory Experiments Measuring Social Preferences Reveal About the Real World?. Introduction, pp. 153–154; Stakes, p. 165; Conclusions, pp. 170–171 (PDF pp. 1–2, 13, 18–19)..
\url{https://bpb-us-w2.wpmucdn.com/voices.uchicago.edu/dist/f/1276/files/2018/09/27-101-25264884-1gx76fi.pdf}.
DOI: 10.1257/jep.21.2.153.
{\small\textit{Earliest verified question/agenda reference; historical priority qualified below. The paper frames generalizability as a central question, identifies scrutiny, context, selection, moral considerations, and stakes as moderators, and calls for theory connecting laboratory and naturally occurring behavior.}}
\end{itemize}
\end{document}
